\documentclass[11pt,a4paper,twoside]{report}

\usepackage[ngerman,USenglish]{babel} % Swap languages with \selectlanguage{...}
\usepackage[utf8]{inputenc} % to write characters (à,Ü,...) in the source
\usepackage[T1]{fontenc} % make characters (à,Ü,...) appear as such in the pdf
\usepackage{lmodern}
\usepackage{amsmath,amssymb,amsfonts} % math stuff
\usepackage{array} % better implementation of tabular and array environments
\usepackage{graphicx} % graphics inclusion
\usepackage{fancyhdr} % beautiful page headers/footers
\usepackage{emptypage} % removes header/footer from empty pages
\usepackage{microtype} % beautiful typesetting
\usepackage{verbatim} % typesetting raw text

% More packages.  Uncomment what you need:

% \usepackage{bbm} % if you need $\mathbbm{1}$
% \usepackage{mathrsfs} % if you need more math script fonts $\mathscr{A}$
% \usepackage{pifont} % more symbols, e.g., circled numbers

\usepackage{pict2e} % a better picture environment
\usepackage{tikz} % versatile drawing facilities
\usetikzlibrary{fit} % used for the example tikzpicture
\usepackage{pgfplots} % versatile plotting facilities
\usepackage{factorGraphs} % package that has macros for factor graphs

% \usepackage{psfrag}
% \usepackage[crop=pdfcrop,cleanup={.tex,.dvi,.ps,.pdf,.log,.aux,.idx,.toc,.out}]
% {pstool} % a pdf-compatible alternative to psfrag.

% \usepackage{amsthm} % or % \usepackage{theorem} % for theorems

% \usepackage{algorithmicx} % typesetting algorithms

% \usepackage[labelfont={md}]{subfig} % for sub-floats, e.g. sub-figures

\usepackage{url} % write URLs

% the last package to load:
\usepackage[bookmarksnumbered=true,hypertexnames=false]{hyperref} % hyperlinks

% Page layout
\setlength{\oddsidemargin}{2.5cm}
\setlength{\evensidemargin}{0.5cm}
\addtolength{\textheight}{0.0cm}
\addtolength{\textwidth}{0.0cm}
\addtolength{\topmargin}{0.0cm}
\setlength{\headheight}{3ex}

\renewcommand{\baselinestretch}{1.2}

% Here the document starts %%%%%%%%%%%%%%%%%%%%%%%%%%%%%%%%%%%%%%%%%%%%%%%%%%%%%%%%%%%%%%%%%%%%%%%%%%%%%%%%%%%%%%%%%%%%%%%%
\begin{document}

\pagenumbering{Roman}
\pagestyle{plain}

% Title page
\begin{titlepage}
  \begin{center}
    \includegraphics[height=15mm]{gfx/ethlogo} \hfill
    \includegraphics[height=15.5mm]{gfx/isilogo_left}
    \rule{\textwidth}{0.5pt}\\[1ex]
    {\Large Spring Semester 2014 \hfill 
      Prof.~H.-A.~Loeliger
    }

    \vspace{\stretch{5}}
    \LARGE Semester Thesis

    \vspace{\stretch{8}}
    \Huge\textbf{
      Phase-Locked Loops
    }
    
    \vspace{\stretch{10}}
    \LARGE{
      Daniel Gilgen\\
      and \\
      Fabio Marti \\
    }
    
    \vspace{\stretch{10}}
  \end{center}
  \rule{\textwidth}{0.5pt}\\[2ex]
  \noindent
  \begin{tabular}{@{}ll@{}}
    \Large Advisor: & \Large Nour Zalmai\\[1ex]
    \Large Co-Advisor: & \Large Lukas Bruderer
  \end{tabular}
\end{titlepage}

% Project description (optional)
\cleardoublepage
(Here the project description may be put in \ldots page 1)
\clearpage
(Here the project description may be put in \ldots page 2)
\cleardoublepage

<<<<<<< HEAD
% Including abstract
% English abstract
\thispagestyle{empty}
\begin{center} \huge \bfseries \abstractname \end{center}
Here comes the abstract \ldots

(The abstract is better included by the command \verb|\include{file}|.)
\clearpage

% Abstract in other language(s) (optional)
\thispagestyle{empty}
\begin{otherlanguage}{ngerman}
  \begin{center} \huge \bfseries \abstractname \end{center}
  Das ist die Zusammenfassung \ldots
  
  (Die Zusammenfassung sollte besser mittels des Befehls \verb|\include{file}|
  eingebunden werden.)
\end{otherlanguage}
\cleardoublepage
=======
% Dedication and/or quotation (optional)
\thispagestyle{empty}
\bigskip
\begin{center}   
  To Arwen
\end{center}

\vfill
\begin{quote}
  \itshape
  ``Among all the conundrums in the world, the Gordian knot is the most
  mysterious.''

  \hfill Anon.
\end{quote}
\vfill

\cleardoublepage

% Acknowledgments
\chapter*{Acknowledgements}
I would like to thank \ldots

(The chapter ``Acknowledments'' is better included by the command
\verb|\include{file}|.)

\vspace{11mm}
\noindent
Zürich, 00 March 2013

\vspace{19mm}
\noindent
Daniel Gilgen \hfill Fabio Marti \hfill \null
\cleardoublepage

% English abstract
\thispagestyle{empty}
\begin{center} \huge \bfseries \abstractname \end{center}
Here comes the abstract \ldots

(The abstract is better included by the command \verb|\include{file}|.)
\clearpage

% Abstract in other language(s) (optional)
\thispagestyle{empty}
\begin{otherlanguage}{ngerman}
  \begin{center} \huge \bfseries \abstractname \end{center}
  Das ist die Zusammenfassung \ldots
  
  (Die Zusammenfassung sollte besser mittels des Befehls \verb|\include{file}|
  eingebunden werden.)
\end{otherlanguage}
\cleardoublepage

% Setup the fancy headers
\pagestyle{fancy}
\renewcommand{\chaptermark}[1]{\markboth{#1}{}}
\renewcommand{\sectionmark}[1]{\markright{\thesection\ #1}}
\fancyhead{}

% Table of contents
\fancyhead[LO]{\scshape \contentsname}
\fancyhead[RO]{\scshape Master's Thesis}
\fancyhead[LE]{\scshape Master's Thesis}
\fancyhead[RE]{\scshape \contentsname}
\pdfbookmark{\contentsname}{contentsbookmarkname}
\tableofcontents
\clearpage

% List of figures (optional)
\fancyhead[LO]{\scshape \listfigurename}
\fancyhead[RE]{\scshape \listfigurename}
\phantomsection\addcontentsline{toc}{chapter}{\numberline{}\listfigurename}
\listoffigures
\clearpage

% List of tables (optional)
\fancyhead[LO]{\scshape \listtablename}
\fancyhead[RE]{\scshape \listtablename}
\phantomsection\addcontentsline{toc}{chapter}{\numberline{}\listtablename}
\listoftables
\cleardoublepage

\setcounter{chapter}{0}
\setcounter{figure}{0}
\pagenumbering{arabic}
\fancyhead[LO]{\rightmark}
\fancyhead[RO]{\scshape \chaptername\ \thechapter}
\fancyhead[LE]{\scshape \chaptername\ \thechapter}
\fancyhead[RE]{\textsc{\leftmark}}

% Chapters %%%%%%%%%%%%%%%%%%%%%%%%%%%%%%%%%%%%%%%%%%%%%%%%%%%%%%%%%%%%%%%%%%%%%%%%%%%%%%%%%%%%%%%%%%%%%%%%

<<<<<<< HEAD
% Chapter 1: Principles of Phase-Locked Loops
\chapter{Principles of Phase-Locked Loops}
% Chapter 1: Principles of Phase-Locked Loops
\section{Introduction}

\section{Blabla}

\clearpage
=======



\begin{figure}[h]
	\centering
    	
	\begin{tikzpicture}
	\factorNode{1}{1}{2}{2}{$A$}
	\factorNode{3}{1}{2}{2}{$B$}
	\factorNode{5}{1}{2}{2}{$A$}
	\rightArrow{7}{7}{2}{$X_k$}
	\leftArrow{9}{6}{2}{$Y_k$}
	\upArrow{1}{4}{2}{$W_k$}
	\downArrow{2}{6}{2}{$V_k$}
	\observable{4}{4}{0.25}{0.25}{$y_k$}
	\forwardMessage{5}{4}{$\hat{W_k}$}
	\forwardMessage{5}{6}{$\hat{W_k}\hat{m_k}$}
	\end{tikzpicture}

  	\caption[Some factor graph elements]
   	{All elements from the file factorGraph.sty.}
	\label{factor_graph}
\end{figure}


\begin{figure}[h]
	\centering
	\begin{tikzpicture}
	\def\x{0}
	\def\y{0}
	\def\s{1.4}
	\def\b{1.4}
	
	\draw node[left] at (\x,\y+\s*4) {$\ldots$};
	\rightArrow{\x}{\y+\s*4}{\s*1}{$X_{k}$}
	\factorNode{\x+\s*1}{\y+\s*4-\b*0.5}{\b}{\b}{$A$}
	\rightArrow{\x+\s*1+\b}{\y+\s*4}{\s}{}
	\factorNode{\x+\s*2+\b}{\y+\s*4-\b*0.5}{\b}{\b}{$=$}
	\rightArrow{\x+\s*2+\b*2}{\y+\s*4}{\s}{$X_{k+1}$}
	\draw node[right] at (\x+\s*3+\b*2,\y+\s*4) {$\ldots$};

	\downArrow{\x+\s*2+\b*1.5}{\y+\s*4-\b*0.5}{\s}{}
	\factorNode{\x+\s*2+\b}{\y+\s*3-\b*1.5}{\b}{\b}{$C$}
	\downArrow{\x+\s*2+\b*1.5}{\y+\s*3-\b*1.5}{\s}{$Y_{k+1}$}
	\factorNode{\x+\s*2+\b}{\y+\s*2-\b*2.5}{\b}{\b}{$+$}
	\downArrow{\x+\s*2+\b*1.5}{\y+\s*2-\b*2.5}{\s}{}
	\observable{\x+\s*2+\b*1.375}{\y+\s-\b*2.75}{\b*0.25}{\b*0.25}{$\tilde{y}_{k+1}$}

	\factorNode{\x+\s*1}{\y+\s*2-\b*2.5}{\b}{\b}{}
	\draw node[below] at (\x+\s*1+\b*0.5,\y+\s*2-\b*2.5) {$\mathcal{N}(0,\sigma^2)$};
	\rightArrow{\x+\s*1+\b}{\y+\s*2-\b*2}{\s}{$Z_{k+1}$}
	\end{tikzpicture}

  	\caption[Exemplary factor graph]
   	{This factor graph has been constructed with custom-built TikZ macros.}
	\label{factor_graph}
\end{figure}

\begin{figure}[h]
	\centering
	\begin{tikzpicture}
	\pllFactorGraph{0}{0}{1}{1}{k-1}{k}
	\end{tikzpicture}

  	\caption[Exemplary factor graph]
   	{This factor graph has been constructed with the factor graph macro from factorGraphs.sty.}
	\label{factor_graph}
\end{figure}





\clearpage


% Chapter 2: Tracking of a Sinusoidal Signal
\chapter{Tracking of a Sinusoidal Signal}
% Chapter 2: Tracking of a Sinusoidal Signal
\chapter{Tracking of a Sinusoidal Signal}

\section{Algorithm}

In \ref{simpleForwardMessagePassing} we can see the most simple approach to construct a PLL with forward message passing.

\begin{algorithm}
\label{simpleForwardMessagePassing}
\begin{algorithmic}[1]
\Parameters $ A, C, n, \sigma^2$
\Require $\{y_k\}_{k=1, \ldots, M}$
\Ensure $\{\tilde{\phi}_k\}_{k=1,\ldots, M}$
\State $k \gets 0$
\State $\overrightarrow{W}_{X_0}\overrightarrow{m}_{X_0} \gets  0^{2n \times 1} $
\State $\overrightarrow{W}_{X_0} \gets  10^{10} \cdot \texttt{identityMatrix}(2n) $

\For{$k < M$}\Comment{Message passing algorithm}
\State $\overrightarrow{W}_{X_{k+1}}\overrightarrow{m}_{X_{k+1}} \gets \texttt{updateMeanVector}(\overrightarrow{W}_{X_k}\overrightarrow{m}_{X_k},y_{k+1},A,C,\sigma^2)$
\State $\overrightarrow{W}_{X_{k+1}} \gets \texttt{updatePrecisionMatrix}(\overrightarrow{W}_{X_k},A,C,\sigma^2)$
\State $\tilde{\phi}_{k+1} \gets \texttt{estimateCurrentPhase}(\overrightarrow{W}_{X_{k+1}}\overrightarrow{m}_{X_{k+1}},\overrightarrow{W}_{X_{k+1}})$
\State $k \gets k+1$
\EndFor

\State \textbf{return} $\{\tilde{\phi}_k\}_{k=1,\ldots, M}$
\end{algorithmic}
\caption[Simple PLL implementation.]
{PLL implementation with forward message passing.}
\end{algorithm}

\clearpage

% Chapter 3: PLL for Harmonic Signals
\chapter{PLL for Harmonic Signals}
% Chapter 3: PLL for Harmonic Signals
\chapter{PLL for Harmonic Signals}

\section{Introduction}

\section{Blabla}

\clearpage

% Chapter 4: BlaBlaBla
\chapter{BlaBlaBla}


% Appendices %%%%%%%%%%%%%%%%%%%%%%%%%%%%%%%%%%%%%%%%%%%%%%%%%%%%%%%%%%%%%%%%%%%%%%%%%%%%%%%%%%%%%%%%%%%%%%%%

% Appendix A
\appendix
\fancyhead[LO]{\rightmark}
\fancyhead[RO]{\scshape\appendixname\ \thechapter}
\fancyhead[LE]{\scshape\appendixname\ \thechapter}
\fancyhead[RE]{\textsc{\leftmark}}

\chapter[Proof of Upper Bound]{Proof for the Lower Bound on the Number of
  Conundrums}
Since we classify the Gordian knot as a conundrum, there is at least one such.

\clearpage

% Bibliography %%%%%%%%%%%%%%%%%%%%%%%%%%%%%%%%%%%%%%%%%%%%%%%%%%%%%%%%%%%%%%%%%%%%%%%%%%%%%%%%%%%%%%%%%%%%%%%%

\fancyhead[LO]{\scshape\bibname}
\fancyhead[RO]{\scshape\appendixname}
\fancyhead[LE]{\scshape\appendixname}
\fancyhead[RE]{\scshape\bibname}

\bibliographystyle{unsrt} % bibliography style in order of first citation

\phantomsection\addcontentsline{toc}{chapter}{\numberline{}\bibname}
\bibliography{IEEEabrv,bibliography}

\end{document}
